\documentclass[12pt,a4paper]{article}

% ============================================================
% PACKAGES
% ============================================================
\usepackage[utf8]{inputenc}
\usepackage[T1]{fontenc}
\usepackage[margin=2.5cm]{geometry}
\usepackage{array}
\usepackage{booktabs}
\usepackage{longtable}
\usepackage[table,dvipsnames]{xcolor}
\usepackage[hidelinks,colorlinks=true]{hyperref}
\usepackage{enumitem}
\usepackage{titlesec}
\usepackage{fancyhdr}
\usepackage[most]{tcolorbox}
\usepackage{tabularx}
\usepackage{setspace}
\usepackage{parskip}
\usepackage{graphicx}

% ============================================================
% COLOR PALETTE
% ============================================================
\definecolor{accentdark}{HTML}{0D4F6B}
\definecolor{accentlight}{HTML}{E6F2F8}
\definecolor{principlebg}{HTML}{FFF6E6}
\definecolor{principleborder}{HTML}{C4900A}
\definecolor{grayrule}{HTML}{CCCCCC}
\definecolor{alertbg}{HTML}{FFF0F0}
\definecolor{alertborder}{HTML}{C0392B}

% ============================================================
% HYPERREF
% ============================================================
\hypersetup{
    linkcolor=accentdark,
    urlcolor=accentdark,
    citecolor=accentdark,
    pdftitle={GenAI-Assisted Document Redaction Governance},
    pdfauthor={Andrea Gennari}
}

% ============================================================
% SECTION STYLING
% ============================================================
\titleformat{\section}
  {\Large\bfseries\color{accentdark}}
  {\thesection}{1em}{}
  [\vspace{2pt}{\color{accentdark}\titlerule[1.2pt]}]
\titlespacing*{\section}{0pt}{24pt plus 4pt minus 2pt}{10pt plus 2pt minus 2pt}

\titleformat{\subsection}
  {\large\bfseries\color{accentdark!85!black}}
  {\thesubsection}{0.8em}{}
\titlespacing*{\subsection}{0pt}{18pt plus 3pt minus 2pt}{6pt plus 2pt minus 1pt}

% ============================================================
% HEADER / FOOTER
% ============================================================
\setlength{\headheight}{16pt}
\addtolength{\topmargin}{-4pt}
\pagestyle{fancy}
\fancyhf{}
\fancyhead[L]{\small\color{accentdark!70!black}\textit{GenAI-Assisted Document Redaction Governance}}
\fancyhead[R]{\small\color{accentdark!70!black}\textit{\nouppercase{\rightmark}}}
\fancyfoot[C]{\small\color{accentdark!70!black}\thepage}
\renewcommand{\headrulewidth}{0.4pt}
\renewcommand{\headrule}{\hbox to\headwidth{\color{accentdark!40!white}\leaders\hrule height \headrulewidth\hfill}}
\renewcommand{\footrulewidth}{0pt}

% ============================================================
% PARAGRAPH & LIST SPACING
% ============================================================
\setlength{\parskip}{6pt plus 1pt minus 1pt}
\setlength{\parindent}{0pt}

\setlist[itemize]{
  topsep=4pt, itemsep=3pt, parsep=1pt,
  leftmargin=1.8em,
  label=\textcolor{accentdark}{\textbullet}
}
\setlist[enumerate]{
  topsep=4pt, itemsep=3pt, parsep=1pt,
  leftmargin=2.2em,
  label=\textcolor{accentdark}{\arabic*.}
}

% ============================================================
% TCOLORBOX STYLES
% ============================================================
\newtcolorbox{definitionbox}[1][]{%
  enhanced, breakable,
  colback=accentlight,
  colframe=accentdark,
  coltitle=white,
  fonttitle=\bfseries,
  boxrule=0.8pt,
  arc=2pt,
  left=8pt, right=8pt, top=6pt, bottom=6pt,
  title={#1},
  attach boxed title to top left={yshift=-2mm, xshift=6mm},
  boxed title style={colback=accentdark, arc=2pt, boxrule=0pt},
}

\newtcolorbox{principlebox}[1][]{%
  enhanced, breakable,
  colback=principlebg,
  colframe=principleborder,
  coltitle=principleborder!80!black,
  fonttitle=\bfseries,
  boxrule=0.8pt,
  arc=2pt,
  left=8pt, right=8pt, top=6pt, bottom=6pt,
  title={#1},
  attach boxed title to top left={yshift=-2mm, xshift=6mm},
  boxed title style={colback=principlebg, arc=2pt, boxrule=0.6pt, colframe=principleborder},
}

\newtcolorbox{alertbox}[1][]{%
  enhanced, breakable,
  colback=alertbg,
  colframe=alertborder,
  coltitle=white,
  fonttitle=\bfseries,
  boxrule=0.8pt,
  arc=2pt,
  left=8pt, right=8pt, top=6pt, bottom=6pt,
  title={#1},
  attach boxed title to top left={yshift=-2mm, xshift=6mm},
  boxed title style={colback=alertborder, arc=2pt, boxrule=0pt},
}

\newtcolorbox{quotebox}{%
  enhanced, breakable,
  colback=accentlight!40!white,
  colframe=accentlight!40!white,
  boxrule=0pt,
  borderline west={2.5pt}{0pt}{accentdark},
  arc=0pt,
  left=10pt, right=8pt, top=6pt, bottom=6pt,
  fontupper=\itshape,
}

% ============================================================
% DOCUMENT BEGIN
% ============================================================
\begin{document}

\begin{titlepage}
    \centering
    \vspace*{3cm}

    {\color{accentdark}\rule{\textwidth}{1.5pt}}

    \vspace{1.2cm}
    {\Huge\bfseries\color{accentdark} GenAI-Assisted Document\par}
    \vspace{0.3cm}
    {\Huge\bfseries\color{accentdark} Redaction Governance\par}
    \vspace{0.8cm}
    {\Large\color{accentdark!80!black} A Governance Framework for Safe Adoption\par}
    \vspace{1.2cm}

    {\color{accentdark}\rule{0.4\textwidth}{0.8pt}}

    \vspace{2cm}
    {\large\color{accentdark!70!black} Submitted by\par}
    \vspace{0.3cm}
    {\LARGE\bfseries\color{accentdark} Andrea Gennari\par}
    \vspace{0.5cm}
    {\normalsize\color{accentdark!70!black} Business and Project Management\par}
    \vfill
    {\color{accentdark}\rule{\textwidth}{1.5pt}}
    \vspace{0.5cm}
    {\large\color{accentdark!70!black} \today\par}
\end{titlepage}

\thispagestyle{empty}
\vspace*{2cm}
\begin{tcolorbox}[
  enhanced,
  colback=accentlight!30!white,
  colframe=accentdark,
  coltitle=white,
  fonttitle=\bfseries\large,
  title=Abstract,
  attach boxed title to top center={yshift=-3mm},
  boxed title style={colback=accentdark, arc=3pt, boxrule=0pt},
  arc=3pt,
  boxrule=0.8pt,
  left=12pt, right=12pt, top=10pt, bottom=10pt,
]
This paper proposes a governance framework for the safe adoption of Generative AI-assisted document redaction tools. The central research question is: how can organizations adopt GenAI-assisted redaction through governance, human-in-the-loop controls, audit logging and risk management? The framework is organized in four pillars anchored to GDPR, the NIST AI Risk Management Framework and ISO/IEC 42001, and is illustrated against the Post Office Horizon data breach case.
\end{tcolorbox}

\newpage

\tableofcontents
\newpage

% ============================================================
\section{Framing and Foundations}
\label{sec:foundations}

\subsection{Framework Overview}

Organizations that publish documents externally need a controlled, auditable process to ensure sensitive content is identified and redacted before release. GenAI tools can support this process, but adopting them without governance introduces new risks: over-reliance on automation, lack of accountability, and potential data leakage.

This paper proposes a \textbf{governance framework} that controls how a document moves from drafting to external publication, structured in four pillars: Governance, Human-in-the-Loop Controls, Audit Logging, and Risk Management. Each pillar addresses a distinct failure mode observed in uncontrolled redaction processes. Table~\ref{tab:structure} shows how the paper is organized.

\begin{table}[h!]
\centering
\small
\rowcolors{2}{accentlight!30!white}{white}
\begin{tabular}{@{} p{4.8cm} p{9cm} @{}}
\toprule
\textbf{\color{accentdark}Section} & \textbf{\color{accentdark}Content} \\
\midrule
\S\ref{sec:foundations}: Foundations & Criteria for safe adoption and regulatory baseline \\
\S\ref{sec:governance}: Pillar~I & Governance, policy and accountability structure \\
\S\ref{sec:hitl}: Pillar~II & Human-in-the-loop controls and document risk classification \\
\S\ref{sec:audit}: Pillar~III & Audit logging and demonstrable compliance \\
\S\ref{sec:risk}: Pillar~IV & Risk management process and register \\
\S\ref{sec:adoption} & Adoption roadmap, KPIs and business case \\
\bottomrule
\end{tabular}
\caption{Structure of the paper.}
\label{tab:structure}
\end{table}

The framework is designed first in full, then deployed through a limited pilot before scaling. This sequencing (predictive design, iterative rollout) follows the hybrid lifecycle approach \cite{falchi_slides}.

\subsection{Defining ``Safe Adoption''}
\label{sec:safe_adoption}

For the framework to be useful it must be possible to verify whether it is working. The following five criteria define what safe adoption looks like in practice; Section~\ref{sec:synthesis} checks the framework against each of them.

\begin{definitionbox}[Five Criteria for Safe Adoption]
\begin{enumerate}
  \item \textbf{No unauthorized disclosure.} No personal or confidential data is shared or published unless it has passed a documented review and approval process. This criterion is directly motivated by the Post Office Horizon incident, where an unredacted document was published.

  \item \textbf{Regulatory compliance.} The adoption satisfies GDPR Art.~25 (data protection by design and by default) and Art.~32 (appropriate technical and organizational measures), and meets the ICO expectation of a systematic, documented pre-publication control.

  \item \textbf{Accountability and auditability.} Every document processed generates a record that can demonstrate to a regulator, a court, or internal governance what action was taken, by whom, and when; this satisfies the GDPR accountability principle (Art.~5(2)).

  \item \textbf{Preserved human oversight.} No high-risk or very-high-risk document is published without explicit, documented human approval. GenAI is used to \emph{support}, not to \emph{replace}, human judgement; the human reviewer is the final control gate.

  \item \textbf{Sustainable adoption.} The framework is genuinely embedded in the workflow: staff are trained, the process is followed consistently, and the KPIs are met within the pilot period. A technically correct framework that is bypassed in practice has zero risk-reduction value.
\end{enumerate}
\end{definitionbox}

Each criterion has a corresponding failure condition that triggers escalation or corrective review (Table~\ref{tab:failure}).

\begin{table}[h!]
\centering
\small
\rowcolors{2}{accentlight!30!white}{white}
\begin{tabular}{@{} p{4.5cm} p{9cm} @{}}
\toprule
\textbf{\color{accentdark}Criterion} & \textbf{\color{accentdark}Failure condition (triggers escalation)} \\
\midrule
No unauthorized disclosure & Any unredacted or confidential content shared/published \\
Regulatory compliance & Any document processed outside the GDPR Art.~25/32 controls \\
Accountability \& auditability & Any document with an incomplete or missing audit record \\
Preserved human oversight & Any high-risk document released without documented human approval \\
Sustainable adoption & KPI targets missed beyond tolerance; systematic bypass of the process \\
\bottomrule
\end{tabular}
\caption{Failure conditions for each safe-adoption criterion.}
\label{tab:failure}
\end{table}

Each criterion maps to one or more KPIs defined in Section~\ref{sec:adoption} and follows the SMART principle: specific, measurable, attainable within the pilot scope, relevant, and time-bound \cite{falchi_slides}.

\subsection{Regulatory and Standards Baseline}
% FILL (next step): GDPR Art.25/32; ICO reprimand expectations; NIST AI RMF (Govern/Map/Measure/Manage); ISO/IEC 27001 & 42001; PMI/PMBOK. End with a mapping table requirement <-> standard <-> where addressed.
\textit{[To be developed in the next step: standards baseline + mapping table.]}

\subsection{Framework Architecture}
% FILL (next step): one-page overview + diagram — governance "spine" supporting the 4 pillars, feeding an adoption lifecycle; restate "AI suggests, humans approve, the organization remains accountable".
\textit{[To be developed in the next step: framework overview + architecture diagram.]}

% ============================================================
\section{Governance (Pillar I)}
\label{sec:governance}
\textit{[To be developed.]}

% ============================================================
\section{Human-in-the-Loop Controls (Pillar II)}
\label{sec:hitl}
\textit{[To be developed.]}

% ============================================================
\section{Audit Logging and Accountability (Pillar III)}
\label{sec:audit}
\textit{[To be developed.]}

% ============================================================
\section{Risk Management (Pillar IV)}
\label{sec:risk}
\textit{[To be developed.]}

% ============================================================
\section{Adoption, Measurement and Business Case}
\label{sec:adoption}
\textit{[To be developed.]}

% ============================================================
\section{Synthesis: Answering RQ1}
\label{sec:synthesis}
\textit{[To be developed.]}

% ============================================================
% PART II — RQ2
% ============================================================
\clearpage
\begin{center}
{\color{accentdark}\rule{\textwidth}{1.2pt}}\\[6pt]
{\Huge\bfseries\color{accentdark} Part II}\\[8pt]
{\LARGE\bfseries\color{accentdark!85!black} RQ2: The Post Office Horizon Case}\\[8pt]
{\color{accentdark}\rule{\textwidth}{1.2pt}}
\end{center}
\vspace{8pt}

\begin{quotebox}
\textbf{RQ2.} How could a GenAI-assisted governance framework have reduced the privacy and data leakage risks in the Post Office Horizon document publication case?
\end{quotebox}

% ============================================================
\section{Case Analysis: Root Causes and Governance Gaps}
\label{sec:rq2_case}
\textit{[To be developed.]}

% ============================================================
\section{Applying the Framework to the Case}
\label{sec:rq2_apply}

\subsection{Governance Controls (Pillar I)}
\textit{[To be developed.]}

\subsection{Human-in-the-Loop Controls (Pillar II)}
\textit{[To be developed.]}

\subsection{Audit Logging (Pillar III)}
\textit{[To be developed.]}

\subsection{Risk Management (Pillar IV)}
\textit{[To be developed.]}

% ============================================================
\section{Counterfactual Analysis}
\label{sec:rq2_counterfactual}

% What would have changed if the framework had been in place?
\textit{[To be developed: for each root cause identified in \S\ref{sec:rq2_case}, show which framework control would have prevented or mitigated it.]}

% ============================================================
\section{Answer to RQ2}
\label{sec:rq2_synthesis}
\textit{[To be developed.]}

% ============================================================
% REFERENCES
% ============================================================
\begin{thebibliography}{9}

\bibitem{falchi_slides}
Alessio Falchi.
\textit{Project \& Business Management --- Course Slides}.
University of Pisa, 2024.

\bibitem{gdpr}
European Union.
\textit{General Data Protection Regulation (GDPR), Articles 5, 25 and 32}.

\bibitem{nist_ai_rmf}
National Institute of Standards and Technology.
\textit{Artificial Intelligence Risk Management Framework (AI RMF 1.0)}.
2023.

\bibitem{iso42001}
International Organization for Standardization.
\textit{ISO/IEC 42001:2023 --- Artificial Intelligence Management Systems}.
2023.

\bibitem{iso27001}
International Organization for Standardization.
\textit{ISO/IEC 27001:2022 --- Information Security Management Systems}.
2022.

\bibitem{ico2025}
Information Commissioner's Office.
\textit{Post Office Limited Reprimand}.
2025.

\end{thebibliography}

\end{document}
